\chapter{LayerNorm 与 RMSNorm}

\begin{introduction}
  \item LN vs RMS
  \item 行统计量
  \item $\mathrm{d}\gamma/\mathrm{d}\beta$ 用 \texttt{atomic\_add}
  \item 推理 API 是 \texttt{rmsnorm}
\end{introduction}

配套：\pyfile{kernel/triton/tutorial/05_norm/}。先跑：\pyfile{step01_layernorm.py} $\to$ \pyfile{step02_online_layernorm.py} $\to$ \pyfile{step03_rmsnorm.py}。推理 Host API 是 \texttt{from step03\_rmsnorm import rmsnorm}。LN 反向写在 step01 同文件。

\section{公式与数据流}

LayerNorm（最后一维 $N$）：
\[
\hat x=\frac{x-\mu}{\sqrt{\sigma^2+\varepsilon}},\qquad
y=\gamma\odot\hat x+\beta.
\]
RMSNorm 没有减均值、没有 $\beta$：
\[
\mathrm{rms}=\sqrt{\mathrm{mean}(x^2)+\varepsilon},\qquad
y=w\odot x/\mathrm{rms}.
\]

Naive LN 把整行 $N$ 一次装下。$N$ 很大时会炸寄存器；online 沿 $N$ 分块累加 $\sum x$、$\sum x^2$。

\fignorm

\src{kernel/triton/tutorial/05_norm/step01_layernorm.py}{69}{81}{LayerNorm 前向：行均值、方差、仿射}

反向里 $\mathrm{d}x$ 每行本地可写；$\mathrm{d}\gamma,\mathrm{d}\beta$ 形状是 $(N,)$，要把所有行加在一起。本块先 \texttt{sum(axis=0)}，再 \texttt{atomic\_add} 到全局。autotune 必须 \texttt{reset\_to\_zero}。

\src{kernel/triton/tutorial/05_norm/step01_layernorm.py}{223}{227}{参数梯度：tile 内求和再原子加}

\section{怎么跑}

\begin{runcmd}
\begin{lstlisting}[style=shell]
python kernel/triton/tutorial/05_norm/step01_layernorm.py
python kernel/triton/tutorial/05_norm/step02_online_layernorm.py
python kernel/triton/tutorial/05_norm/step03_rmsnorm.py
\end{lstlisting}
\end{runcmd}

\section{正确性}

\begin{note}
LN 对 \texttt{F.layer\_norm}。RMS 对 \texttt{F.rms\_norm}。fp32，$\mathrm{rtol}$ 约 $10^{-4}$。训练必须 \texttt{TritonLayerNorm.apply}，不能对裸 \texttt{layernorm(...)} 求 \texttt{grad}。
\end{note}

\section{效果}

\begin{remark}
大 $N$ 时分块版接近 torch，是因为 naive 整行会炸寄存器，不是因为公式更快。教学 LN 大 hidden 仍打不过 \texttt{F.layer\_norm}（背后是 cuDNN），训练默认不换它。
\end{remark}
